% SPDX-License-Identifier: CC-BY-SA-4.0
% Session: Form parallelism to concurrency
% Exercise: The mad scientist

\begingroup

\begin{exercice}[Le savant fou%
    \footnote{Inspiration : M. Herlihy \& N. Shavit. \emph{The Art of Multiprocessor Programming}. Morgan Kaufmann (2008)}]
  \label{exo:introduction/mad_scientist}

  En visite pédagogique dans un laboratoire, vous et le reste de votre groupe êtes 
  capturés par un savant fou. Il fait l'annonce suivante :

  \og
  J'ai mis en place une pièce qui contient une lumière et un interrupteur
  qui la commande. Actuellement, la lumière est éteinte.
  Vous pouvez discuter ce soir d'une stratégie commune, mais à partir de demain
  vous serez isolés et ne pourrez plus communiquer.
  Chaque jour, je tirerai au hasard de manière uniforme l'un d'entre vous, qui passera la journée dans la pièce.
  À tout moment durant sa visite, il pourra déclarer~: \emph{Nous avons tous visité la pièce}.
  S'il a raison, vous êtes libres. Sinon, je vous reboote les cerveaux !
  \fg

  Hypothèses implicites :
  \begin{itemize}
  \item La lumière est éteinte au départ.
  \item L'interrupteur est binaire (allumé/éteint).
  \item La lumière est invisible depuis l'extérieur de la pièce,
    et il n'y a aucun autre moyen de communication entre les étudiants.
  \end{itemize}
  
  \begin{question}
  \item Énoncez précisément ce que doit garantir une stratégie réussie.
  \end{question}
  \begin{correction}
    Une stratégie correcte doit vérifier :
    \begin{description}
    \item[Sûreté.] Si un étudiant déclare \og Nous avons tous visité la pièce \fg,
      alors tous les étudiants l'ont effectivement visitée.
    \item[Vivacité.] Avec probabilité 1, un étudiant finira par pouvoir faire cette déclaration.
    \end{description}
  \end{correction}
  
  \begin{question}
  \item Proposez une stratégie et justifiez brièvement en quoi elle satisfait l'exigence précédente.
  \end{question}
  \begin{correction}
    Une stratégie possible consiste à désigner un étudiant comme \emph{délégué} :
    \begin{itemize}
    \item Chaque étudiant non délégué allume la lumière la première fois qu'il entre
      dans la pièce et la trouve éteinte, puis ne la rallume plus jamais.
    \item Chaque fois que le délégué entre dans la pièce et voit la lumière allumée,
      il l'éteint et incrémente un compteur local.
    \item Lorsque son compteur atteint $n-1$, il déclare que tous les étudiants ont visité la pièce.
    \end{itemize}
    Cette stratégie garantit la sûreté (le délégué ne peut compter qu'un allumage par étudiant)
    et la vivacité presque sûrement (le tirage aléatoire finit par sélectionner chaque étudiant
    au moins une fois alors que la lumière est éteinte).
  \end{correction}
  
  \begin{question}
  \item Reprenez la question précédente, sans l'hypothèse sur l'état initial de la lumière.
  \end{question}
  \begin{correction}
    Si l'état initial de la lumière est inconnu,
    le protocole peut être adapté en demandant au délégué d'ignorer le premier allumage observé
    (pour éliminer l'incertitude sur l'état de départ).
  \end{correction}
  
  \begin{question}
  \item Reprenez les deux premières questions, sans l'hypothèse que le tirage est uniforme. À la place, on suppose
    que si personne ne déclare que tout le monde a visité la pièce, chaque étudiant est sélectionné une infinité de fois
    (pas nécessairement à intervalles bornés). 
  \end{question}
  \begin{correction}
    Si le tirage n'est pas uniforme mais seulement équitable
    (chaque étudiant est choisi une infinité de fois, sans borne sur l'intervalle entre deux visites),
    la stratégie reste correcte : la sûreté n'est pas affectée,
    et la vivacité est renforcée et devient déterministe : un étudiant finira par pouvoir faire cette déclaration.
  \end{correction}

\end{exercice}

\endgroup
\endinput
